\documentclass[12pt]{article}
\usepackage{amsmath}
\usepackage[margin=1in]{geometry}

\title{Labwork 2: Linear Regression}
\author{Trung Tran Minh - 2540024}
\date{}

\begin{document}
\maketitle

\section{Introduction}

In this lab, I implement linear regression from scratch by using gradient descent.
The input is a CSV file with two columns. The first column is $x$ and the second column is $y$.
The model is:

\[
\hat{y} = w_1x + w_0
\]

The output of the program is $w_1$ and $w_0$. The dataset used in this lab is $lr.csv$

\section{Result}

I use learning rate: r = 4e-3


\section{Effect of Different Learning Rates}

I also try different learning rates. The result is:

\[
\begin{array}{c|c|c|c|c}
r & w_1 & w_0 & loss & comment \\
\hline
0.00001 & 1.8969 & 11.4190 & 381.4041 & slow \\
0.0001 & 1.3176 & 44.7973 & 21.6618 & better \\
0.0004 & 1.2622 & 47.9869 & 18.9024 & good \\
0.001 & -28.8146 & -0.4482 & 2313896.1725 & diverge
\end{array}
\]

From this table, when the learning rate is too small, the model learns very slowly.
When the learning rate is $0.0004$, the model converges well.
When the learning rate is too large, like $0.001$, the loss becomes very large and the algorithm diverges.

\end{document}
